\section{Exercise 1}

Mảng 5 phần tử khai báo bằng \texttt{.word 0:5}. \texttt{\$s0} giữ địa chỉ đầu mảng, \texttt{\$t0} là chỉ số phần tử đang xử lý:

\begin{codeblock}[language=MIPS]
  array: .word 0:5

  la $s0, array    # địa chỉ đầu mảng
  li $t0, 0        # chỉ số index
\end{codeblock}

Mỗi vòng lặp đọc một số nguyên, chuyển số đó về bội số của 4 gần nhất rồi lưu vào mảng.

Phần dư lấy bằng \texttt{andi \$t1, \$v0, 3}. Chia cho $4 = 2^2$ là bỏ đi hai bit thấp, nên phần dư nằm gọn trong hai bit đó; AND với $3 = 11_2$ giữ lại đúng hai bit ấy và xoá phần còn lại. Cách này chỉ cần một lệnh, không dùng đến \texttt{HI}/\texttt{LO} như \texttt{div}, và vì là bitwise nên số âm cũng ra phần dư $0$–$3$ như phép chia dư không âm.

Phần dư 0 thì giữ nguyên; dư 1 hoặc 2 thì trừ đi phần dư (làm tròn xuống); dư 3 thì cộng thêm 1 (làm tròn lên). Vì thế 35 thành 36, còn 34 và 33 thành 32.

\begin{codeblock}[language=MIPS]
process_loop:
  li $v0, 5                # syscall 5: đọc số nguyên
  syscall

  andi $t1, $v0, 3         # $t1 = n % 4
  sub $v0, $v0, $t1        # n -= n % 4: làm tròn xuống
  beq $t1, 3, round_up     # dư 3 thì bội số gần nhất nằm phía trên
  j next

round_up:
  addi $v0, $v0, 4         # n += 4: nhảy lên bội số kế tiếp

next:
  sll $t1, $t0, 2          # $t1 = 4i
  add $t1, $t1, $s0        # $t1 = base + 4i
  sw $v0, 0($t1)           # lưu vào phần tử thứ i
  addi $t0, $t0, 1         # i++
  bne $t0, 5, process_loop # đủ 5 phần tử thì dừng
\end{codeblock}

Vòng lặp thứ hai đọc lại mảng và in từng phần tử, mỗi số cách nhau một dấu cách:

\begin{codeblock}[language=MIPS]
  li $t0, 0         # i = 0: duyệt lại từ đầu mảng
print:
  li $v0, 1         # syscall 1: in số nguyên
  sll $a0, $t0, 2   # $a0 = 4i
  add $a0, $s0, $a0 # $a0 = base + 4i
  lw $a0, 0($a0)    # $a0 = phần tử thứ i
  syscall

  li $v0, 11        # syscall 11: in ký tự
  li $a0, ' '       # in dấu cách
  syscall

  addi $t0, $t0, 1  # i++
  bne $t0, 5, print
\end{codeblock}

Chương trình không in lời nhắc khi đọc, nên nhập lần lượt 35, 34, 33, 12, 7 thì in ra:

\begin{codeblock}[language={}]
36 32 32 12 8
\end{codeblock}
